% WS4 report fragment (issue #42): system and cluster design (Section 6), the core run and the E5
% results and analysis (Sections 8 and 9), limitations, scope changes and conclusion (Section 10) and
% the WS4 contribution (Section 11). No preamble; docs/report/main.tex pulls it in with \input.
% Every number carries a "% source:" comment; "% TODO" marks a result that does not exist yet.
% Run time against cores and the second machine (E8) is reported in the WS1 fragment, not repeated here.
% Check its length alone with: cd docs/report/standalone && tectonic ws4.tex

\newcommand{\wsfourSystem}{%
\paragraph{Data flow.} The Stack Overflow archive is read as a stream and never unpacked; the ETL
writes posts without text, question text (\texttt{python} questions only) and post links as Parquet,
partitioned by year. The pipeline then runs in Spark~\cite{zaharia2016spark} as four stages:
shingles from question titles, MinHash signatures, LSH candidate pairs, and evaluation against the
answer key. Each stage writes its output to HDFS, so a later stage can be repeated without the
earlier ones. Evaluation appends one row of metrics per run and writes a table with one row per
labelled pair for the error analysis.
% source: src/etl/README.md; scripts/cluster/stages.py; docs/CONTRACTS.md

\paragraph{Cluster.} Two laptops with 16~GB of memory each form a Spark 3.5.3 standalone cluster
with Hadoop 3.3.6, joined by a wired Gigabit link. Laptop~A (Intel i7-13650HX) runs the Spark
master, the HDFS NameNode, a worker and a DataNode; Laptop~B (AMD Ryzen~7 5800H) runs a worker and a
DataNode. HDFS keeps two copies of each block. In the full core run Laptop~A gave 8 cores and 8~GB
to its worker and Laptop~B 4 cores and 5~GB, with 4~GB per executor and 200 shuffle partitions.
% source: results/core_pipeline_summary.md
The scaling runs changed the number of cores as described with their results.
% source: docs/e8_runbook.md

\paragraph{Operation.} Cluster configuration files are generated from templates for each laptop, and
scripts start, stop, check and smoke-test the services. Spark event logs go to HDFS and are the
source of run times, stage times, shuffle volume and task times. Pipeline settings are in one JSON
file, and seeds are fixed.
% source: docs/cluster_setup.md; configs/full_run.json
}

\newcommand{\wsfourResults}{%
\subsection{Core Run and Baselines (E5)}
\paragraph{Core run.} Table~\ref{tab:ws4-core} gives both LSH settings on 10\% and 100\% of the
\texttt{python} subset. At full size, building shingles and signatures took 1,065~s and wrote
18.1~GB of shuffle data; candidate generation and evaluation with the chosen setting ($b=32$, $r=4$)
took 3,571~s and wrote 105.6~GB, with at most 976~MB of execution memory in one task.
% source: results/core_pipeline_summary.md

\begin{table}[h]
  \centering
  \small
  \setlength{\tabcolsep}{3pt}
  \caption{Core runs. Each size is scored on the labelled pairs inside it.}
  \label{tab:ws4-core}
  \begin{tabular}{@{}llrrr@{}}
    \toprule
    Data & $b$, $r$ & Candidates & Recall@10 & MAP \\
    \midrule
    10\% & 16, 8 & 121,120 & 0.008 & 0.007 \\ % source: results/metrics.csv
    10\% & 32, 4 & 27,186,575 & 0.064 & 0.042 \\ % source: results/metrics.csv
    100\% & 16, 8 & 11,725,463 & 0.012 & 0.009 \\ % source: results/metrics.csv
    100\% & 32, 4 & 740,022,352 & 0.038 & 0.023 \\ % source: results/metrics.csv
    \bottomrule
  \end{tabular}
\end{table}

\paragraph{Exact search.} Two exact all-pairs implementations (an inverted-index join and a direct
Spark join) returned the same pairs with Jaccard at least 0.1 on nested Super User samples:
4,883,441 pairs for 10,000 questions and 98,759,691 for 50,000. In the recorded runs the direct join
was slower at 10,000 (40~s against 33~s) and faster at 50,000 (355~s against 449~s). Those timings
include counting and writing an uncached result on Windows, so they show only rough cost.
% source: results/e5_baselines/runs.csv; results/e5_baselines/README.md

\paragraph{Shared sample.} Table~\ref{tab:ws4-shared} runs exact search, our LSH and Spark's
\texttt{MinHashLSH} on the same 9,650 Super User questions with 2,139 labelled pairs. Run times are
left out: two runs of the same command differed by up to 4.6 times.
% source: results/e5_baselines/ours_vs_spark.md

\begin{table}[h]
  \centering
  \small
  \setlength{\tabcolsep}{3pt}
  \caption{One shared Super User sample. The last two columns are the share of exact pairs with at
  least that similarity that the method also returns.}
  \label{tab:ws4-shared}
  \begin{tabular}{@{}lrrrr@{}}
    \toprule
    Method & Candidates & R@10 & $J\ge0.4$ & $J\ge0.6$ \\
    \midrule
    Exact, $J\ge0.1$ & 3,671,388 & 0.351 & 100\% & 100\% \\ % source: results/e5_baselines/shared_sample.csv
    Ours, 64, 2 & 3,660,207 & 0.338 & 100\% & 100\% \\ % source: results/e5_baselines/shared_sample.csv
    Ours, 32, 4 & 78,091 & 0.116 & 69\% & 100\% \\ % source: results/e5_baselines/shared_sample.csv
    Ours, 16, 8 & 146 & 0.024 & 5\% & 62\% \\ % source: results/e5_baselines/shared_sample.csv
    Spark, 5 tables & 99 & 0.024 & 5\% & 100\% \\ % source: results/e5_baselines/shared_sample.csv
    \bottomrule
  \end{tabular}
\end{table}
}

\newcommand{\wsfourAnalysis}{%
\subsection{What the Baselines and the Cluster Show}
\paragraph{Against exact search.} Exact search is itself limited: it reaches Recall@10 0.351, and
only 68.5\% of the labelled pairs share enough title words to pass $J\ge0.1$. With $b=32$, $r=4$ our
LSH returns 47 times fewer candidates and reaches about a third of that recall; $b=64$, $r=2$ comes
close to it but removes almost no candidates. Spark's hash tables are combined with OR only, and its
join then keeps only the pairs inside a distance cutoff, here $J\ge0.6$. Its low recall therefore comes
mainly from the cutoff, and the comparison does not show that either LSH is better at equal settings.
% source: results/e5_baselines/ours_vs_spark.md

\paragraph{Reading the core runs.} Recall@10 is higher on 10\% (0.064) than on all of the subset
(0.038). The two rows are scored on different labelled pairs, and at full size each pair competes
with more candidates for ten places while more buckets pass the size limit. We did not separate
these effects.
% source: results/metrics.csv

\paragraph{Running on two laptops.} Run times on more cores and on the second machine are analysed
with the speedup results. One run was lost when an automatic system update replaced Java during the
job; the services were restarted and that test's runs were done again. In every setting of the speedup tests the three run times stayed within 13.8\% of their median. The effect of data size on run time and memory
has not been measured.
% source: docs/e8_runbook.md
% TODO(#43): add the data-size results, or state that the test was not run.
}

\newcommand{\wsfourConclusion}{%
\paragraph{Limitations.} Shingles come from titles only, and most labelled duplicates share few title
words. The duplicate labels are incomplete. Full-scale results cover only \texttt{python} questions,
each full-size run time comes from a single run, and the baseline comparison uses Super User. The recall lost
only to the bucket size limit was not measured.

\paragraph{Scope changes.} Our first proposal also covered tag experts (PageRank), stream counting
and the effect of AI tools on Stack Overflow. We dropped these to study one problem in depth, and
kept SON as a secondary question about blocking. Splitting big buckets was implemented but not run at
full size. The exact search was not run on 100,000 questions.
% source: docs/proposal/ENSF612_Project_Proposal.tex; docs/proposal/ENSF612_Project_Proposal_v3.tex
% TODO(#43): state what happened to the data-size test and the run on all of Stack Overflow.

\paragraph{Conclusion.} \emph{Accuracy:} on the full \texttt{python} subset our chosen setting
reaches Recall@10 3.8\% while ruling out 99.97\% of all pairs; on the Super User sample it reaches
about a third of exact search's recall.
% source: results/metrics.csv; results/e5_baselines/shared_sample.csv
\emph{Theory:} the spread of our MinHash estimates is within about 10\% of the formula, and measured
LSH curves differ from the theory by 1 to 3 points on average (Super User and a Stack Overflow sample).
% source: results/e2_signature/README.md; results/e3_lsh_real/README.md; results/e3_lsh_stackoverflow/tradeoff.csv
\emph{Scale:} without a limit, 659 of 46.4 million buckets would hold 89.9\% of all bucket pairs;
on a quarter of the subset, 8 cores on one laptop gave a speedup of 3.47, and two laptops reading HDFS were 2.01 times as fast as
the Intel laptop alone. How run time grows with data size is still open.
% source: results/e7_buckets/bucket_profile.csv; results/e8_scaling/analysis/README.md
\emph{Blocking:} SON groups cut candidates by 60\% but run time by only 8\%, and Recall@10 fell from
0.038 to 0.027, because 54\% of labelled pairs fall into different groups.
% source: results/e6_blocking/blocking_comparison.csv
The main obstacle is the text: with title words even exact search over all pairs with $J\ge0.1$ stops at Recall@10 0.351 on the
Super User sample, so much better recall needs text in which duplicates overlap more, which we have
not tested at full size.
}

\newcommand{\wsfourContribution}{%
\paragraph{Md Jubayer Khan (WS4: Cluster, Baselines and Scale).} Set up the Spark and HDFS cluster
on Laptop~A with its runbook; added the archive check and the full data-loading wrapper for Stack
Overflow; implemented the exact all-pairs baseline and the Spark \texttt{MinHashLSH} baseline (E5);
recorded the full core pipeline run; and produced the Stack Overflow dataset profile.
% TODO(team): Md Jubayer Khan to confirm this paragraph.
}
